\section{Related work}
\label{sec:related-work}

The protocol touches five bodies of work: the definition and taxonomy of leakage, the bias of resampling estimates when a data-dependent step sits outside the resampling loop, the validation of models whose data have a time order, the practice of credit scoring and of the validation tests of a bank, and the software that tests models and data. Each is characterised here by what it assumes, what it gives and where it stops short of a protocol that a validator can run per case.

\subsection{Definitions and taxonomies of leakage}

Kaufman, Rosset and Perlich formulate leakage in data mining as the introduction of information about the target that should not legitimately be available for learning, give a methodology for detecting and avoiding it, and illustrate it with cases from data-mining competitions \citep{kaufman2012}. Kapoor and Narayanan survey the consequences in machine-learning-based science and give a taxonomy of the ways in which a training set and a test set come to share information, from the preprocessing of the pooled data to temporal leakage \citep{kapoor2023}. Both works define the failure and list its forms; neither fixes a procedure by which a particular validation is shown to be free of it. The adaptive-analysis literature supplies the neighbouring result that a holdout loses its validity when it is reused to steer the analysis, so that the number measured on it is optimistic by an amount that grows with the number of adaptive choices \citep{dwork2015,blum2015}, and that new test sets drawn after the benchmark was fixed give lower accuracy than the reused one \citep{recht2019}. The protocol of this paper borrows the notion of an information set from these works and turns it into a property that can be tested on a fitted pipeline: the parameters do not change when the outside information changes.

\subsection{Preprocessing and selection inside the resampling loop}

The cleanest statement of the problem for supervised steps is the one by Ambroise and McLachlan: a feature selection that uses the labels of all samples before cross-validation gives an error rate that is far too optimistic, and the selection has to be repeated inside each training fold \citep{ambroise2002}. Hastie, Tibshirani and Friedman present the same point as the wrong and the right way to cross-validate \citep{hastie2009}. Varma and Simon show that the error of the cross-validated best model is biased downwards when the same folds choose the model and report its error, and that a nested design removes the bias \citep{varma2006}; Cawley and Talbot reach the same conclusion for the tuning of hyperparameters \citep{cawley2010}. Moscovich and Rosset analyse the bias that even unsupervised preprocessing introduces when it is applied before the folds \citep{moscovich2022}, and Vabalas and colleagues measure how the optimism of the resampling estimate changes with the sample size for several validation designs \citep{vabalas2019}. Arlot and Celisse survey the cross-validation procedures \citep{arlot2010}, and Bates, Hastie and Tibshirani show that the cross-validated error estimates the average error of models fitted on new training sets rather than the error of the one model that was fitted \citep{bates2023}. These works give the direction of the bias for resampling and, in the cases they measure, its size on one data type. They do not give a gate that can be attached to a production pipeline with its own preprocessing, and they do not treat the calibration map as a fitted step with its own slice.

\subsection{Validation across time and structure}

Where the data have a time order, a random split inside the sample reports a performance that the future need not share. Bergmeir and Ben\'{\i}tez evaluate cross-validation for time series predictors and compare it with blocked and out-of-sample designs \citep{bergmeir2012}; Bergmeir, Hyndman and Koo show that the usual cross-validation is valid for purely autoregressive models with uncorrelated errors and not beyond them \citep{bergmeir2018}. Roberts and colleagues review the blocking strategies that data with temporal, spatial, hierarchical or phylogenetic structure require \citep{roberts2017}, and the concept-drift literature explains why the relation between inputs and outcome can change after the sample was drawn \citep{gama2014}. In credit risk the time order is the rule: the outcome of a loan is observed years after the observation date, so that a training window can be defined only by the vintages whose outcome is already closed at the date on which the model is built. The protocol turns that requirement into an explicit rule and measures, on five cohorts of a public loan file, how far the random split of the window is from the declared out-of-time cohort.

\subsection{Credit scoring, calibration and the tests of a validation function}

Scorecards are built with supervised binning and the weight of evidence, selected by information value and fitted by logistic regression \citep{siddiqi2017,thomas2002,hand1997,navas2020}. Benchmarks compare them with learned challengers \citep{baesens2003,lessmann2015,gunnarsson2021}, and the data set of Yeh and Lien, with 30,000 card holders, is a common testbed \citep{yeh2009}. Calibration maps are fitted on data that the model did not train on, by logistic or isotonic regression \citep{niculescu2005,zadrozny2002}, and corrections for class imbalance change the calibration of the resulting probabilities \citep{goorbergh2022}. The tests of a validation function are the discrimination of the current sample against its development value with the variance of the AUC \citep{hanley1982,delong1988,ecb2019,bcbs14}, the comparison of nested models, whose DeLong test has known misuses \citep{demler2012}, and the stability indices, whose distribution is known only approximately \citep{yurdakul2020}. None of them is designed to detect that a number was inflated by the fitting path, and the measurement in Section~\ref{sec:results} shows what they see of one kind of leak.

\subsection{Software that tests models and data}

Open toolkits ship tests that look for leakage in data and models. The library \texttt{mercury-robust} of BBVA provides a label-leaking test and a sample-leaking test among its data tests \citep{mercury2023}, and \texttt{probatus} of ING provides a resemblance model between two samples, which is the test of how well a classifier separates development rows from out-of-time rows \citep{probatus2026}. A reading of the source of the interpretable-classifier project \texttt{auto-bayesian} shows the practice that the protocol is meant to exclude: the candidate structure is selected by validation AUC and the reported metrics are computed on the same validation partition \citep{autobayesian2026}. Metamorphic testing is the software-engineering technique closest to the gates: a relation between two executions of a program on related inputs is checked instead of an expected output \citep{chen2018,segura2016}, and the rubric of production readiness for machine-learning systems lists tests on data and on training among its items \citep{breck2017}. The three gates of this paper are metamorphic relations of the fitting path: if the outside rows are changed, every fitted parameter must be equal.

\subsection{Position of the proposed protocol}

The protocol stands between the taxonomies, which say what leakage is, and the toolkits, which look for symptoms of it in data and in scores. It states the row set that each step of one named pipeline may read, fixes the split per case, and tests the pipeline by perturbation instead of by the symptom, so that a leak through a step is detected whatever its effect on the score. The measurement that accompanies it, the cost of each leak in bias and coverage and the power of each gate, is what a validator needs to decide how much weight a holdout number can bear. A validator can adopt the protocol without changing the model: the gates only run the fitting path again on perturbed copies of the data, and the measurement says how large the consequence of each failure has been on data of comparable shape. The statistical design follows the recommendations for simulation studies \citep{morris2019}, with Wilson intervals for proportions \citep{brown2001} and the sub-Gaussian bounds of Hoeffding and of the concentration literature for the theory \citep{hoeffding1963,boucheron2013}, and the pipeline is implemented with scikit-learn \citep{pedregosa2011}.
